\section{A Token at the Router}
\label{sec:ch15-a-token-at-the-router}
\index{Cosmo Router!JWT validation|(}%

The router reads tokens because a block in \code{router/config.yaml} tells it
how. Mosaic's is three lines of substance:

\begin{minted}[fontsize=\footnotesize,breaklines]{yaml}
authentication:
  jwt:
    jwks:
      - symmetric_algorithm: HS256
        secret: "${MOSAIC_JWT_SECRET}"
        header_key_id: mosaic-dev
\end{minted}

\index{authentication@\code{authentication}!the symmetric branch}%
The third line is not optional and I found that out the way you would expect.
The schema for this block is a \code{oneOf}: one branch wants a \code{url} and
the other wants \code{secret}, \code{symmetric\_algorithm} and
\code{header\_key\_id} together. Leave the key id out and the container does not
start:

\begin{minted}[fontsize=\scriptsize,breaklines,breakanywhere]{text}
Could not load config: errors while loading config files: router config
validation error for /etc/mosaic-router/config.yaml: jsonschema validation failed
with 'https://raw.githubusercontent.com/wundergraph/cosmo/main/router/pkg/config/config.schema.json#'
- at '/authentication/jwt/jwks/0': oneOf failed, none matched
  - at '/authentication/jwt/jwks/0': validation failed
    - at '/authentication/jwt/jwks/0': missing property 'url'
    - at '/authentication/jwt/jwks/0': not failed
  - at '/authentication/jwt/jwks/0': missing property 'header_key_id'
\end{minted}

Read that against chapter~\ref{ch:real-time-in-a-federated-world}, where wgc
took a misspelled key without a word and quietly used a default. Two tools in
one ecosystem, and they disagree about whether an unknown or missing
configuration key is a bug report or a shrug. The router is right and it is not
close: this message names the file, the JSON pointer and both branches it tried.

\index{header\_key\_id@\code{header\_key\_id}!required to start, not to validate}%
Having insisted on a key id, the router then does not require one of a token. I
minted one with no \code{kid} header at all, which
\code{scripts/mint-token.mjs --no-kid} still does, and the router accepted it
and answered the query. The configuration is mandatory and the thing it names is
not checked. That is a small finding and it is the kind that costs an afternoon,
because the natural conclusion from the startup failure is that key ids matter
here.

\subsection*{What it logs, which is nothing}

\index{Cosmo Router!silence about authentication}%
Start the router with that block and compare its output to a run without it.
There is no line saying JWT validation is on. Not at \code{info}, and not at
\code{debug}: grepping the whole startup log for \code{auth}, \code{jwt} and
\code{jwk} matches nothing at all.

That is worth sitting with for a second, because chapter~\ref{ch:enter-the-router}
made a point of how much this process says about itself at startup. It names its
listen address, its playground, its introspection setting, the config file it is
watching, and every Cosmo Cloud feature it is doing without. It does not mention
that it has begun checking signatures.

The two ways to find out are both behavioural. Send a token that is wrong and
you get a \code{401}. Send one that is right and the response carries an
\code{X-Authenticated-By: jwks} header. Neither of those is something you would
think to look for if you were trying to answer \enquote{did my config load}, and
the first thing most people do with a security setting is check that it took.

\subsection*{Anonymous, wrong, and expired are three different answers}

\index{authorization@\code{authorization}!the block and its defaults}%
The companion block decides what happens to a request with no token at all:

\begin{minted}[fontsize=\footnotesize,breaklines]{yaml}
authorization:
  require_authentication: false
  reject_operation_if_unauthorized: false
  enable_pre_fetch_field_authorization: true
\end{minted}

The first two are defaults written out, because both are decisions. The first of
them, left false, is what keeps the storefront public, and Mosaic needs that: it
sells to people who have not signed in. A field that needs a caller says so
itself, and an anonymous request is refused that field whatever this setting is.
The global switch decides whether an anonymous request may enter the graph at
all, not whether it may have everything in it. Which layer does the refusing is
a separate question with a more interesting answer, and
section~\ref{sec:ch15-a-connection-with-no-identity} is where I measure it.

So there are three outcomes rather than two. No token gets you in as nobody. A
token that will not validate, including an expired one, gets a flat \code{401}
before anything is planned:

\begin{minted}[fontsize=\footnotesize,breaklines]{text}
HTTP/1.1 401 Unauthorized
{"errors":[{"message":"unauthorized"}]}
\end{minted}

\index{tokens!ClockSkew@\code{ClockSkew}}%
A token expires when it says it does, in this repository, because the shared
security defaults set \code{ClockSkew} to zero. The default is five minutes,
which is the right production setting and makes a one-second token useless as a
demonstration.

And a token that validates but lacks the scope a field wants gets refused with
the scope named back:

\begin{minted}[fontsize=\scriptsize,breaklines]{json}
{"errors":[{"message":"Unauthorized to load field 'Query.customerById', Reason:
  missing required scopes.","path":["customerById"],
  "extensions":{"code":"UNAUTHORIZED_FIELD_OR_TYPE"}}],
 "data":{"customerById":null},
 "extensions":{"authorization":{"missingScopes":[{"coordinate":
   {"typeName":"Query","fieldName":"customerById"},"required":[["customers:read"]]}],
  "actualScopes":["reviews:write"]}}}
\end{minted}

\code{required} is a list of lists, and section~\ref{sec:ch15-two-vocabularies}
is where that shape comes from. What I want to point at here is that the router
tells a client exactly which scope it was missing and exactly which ones it
holds. That is generous, and whether it is too generous is a question
chapter~\ref{ch:security-hardening} should answer rather than this one.

\subsection*{The setting that decides whether a refusal costs anything}

\index{authorization!pre-fetch field authorization}%
The third line is the one that is not a default, and it is there because of a
number rather than a principle.

Left off, which is how the router ships, an anonymous request for
\code{Query.customerById} is planned, sent to Accounts, answered, and then the
field is discarded. Accounts logs one \code{Mosaic.RequestTimeline} line per
request it serves, so count those either side of one refused request:

\begin{minted}[fontsize=\footnotesize,breaklines]{text}
enable_pre_fetch_field_authorization   requests Accounts served
off (the default)                      1
on                                     0
\end{minted}

The client can see it too, because with the setting off the answer carries both
refusals, the subgraph's nested inside the router's:

\begin{minted}[fontsize=\scriptsize,breaklines]{json}
{"errors":[{"message":"Failed to fetch from Subgraph 'accounts'.","extensions":
  {"errors":[{"message":"The current user is not authorized to access this resource.",
   "path":["customerById"],"extensions":{"code":"AUTH_NOT_AUTHENTICATED"}}],
   "serviceName":"accounts","statusCode":200}},
  {"message":"Unauthorized to load field 'Query.customerById', Reason: not
   authenticated.","path":["customerById"],
   "extensions":{"code":"UNAUTHORIZED_FIELD_OR_TYPE"}}],
 "data":{"customerById":null}}
\end{minted}

and with it on, one:

\begin{minted}[fontsize=\scriptsize,breaklines]{json}
{"errors":[{"message":"Unauthorized to load field 'Query.customerById', Reason: not
   authenticated.","path":["customerById"],
   "extensions":{"code":"UNAUTHORIZED_FIELD_OR_TYPE"}}],
 "data":{"customerById":null}}
\end{minted}

One round trip is not worth a configuration change. What is worth it is the
other half: with the setting off, data the caller was never allowed to see was
loaded out of a database and assembled in a process on their behalf, and then
thrown away. Nothing went wrong. Something has to go wrong later, in how a
response is put together or logged or traced, for that to matter, and the point
of not loading it is that the later thing then has nothing to leak.

Both rows are gated, in two different ways, because they are visible in two
different places. Both verification scripts count the Accounts requests a
refused field costs and require zero. And a case in
\code{scripts/router-cases.mjs} starts two routers differing only in that
setting and asserts that the nested \code{serviceName} error is present in one
reply and absent from the other, which is the same fact as the count without
needing a subgraph's console to read it.

One caution, which took me until
section~\ref{sec:ch15-a-connection-with-no-identity} to measure and which I
would rather you carried from here: the zero in that table is a root field being
refused, and a root field is a whole fetch. A guarded field hanging off an
entity that the caller is also entitled to read is a different case, and the
setting does less there than its name suggests.

\medskip
That is the router holding up its end. What it is enforcing, though, came out of
the subgraphs' own schemas, and the way those two families of annotation reach
it is stranger than it looks.

\index{Cosmo Router!JWT validation|)}%
